\chapter*{Abstract}
\addcontentsline{toc}{chapter}{Abstract}
In maritime radar, sea clutter is the unwanted return from the sea surface and is a form of interference that affects target detection. This clutter is spiky and well modelled as K-distributed. Commonly used adaptive detectors assume Gaussian-modelled clutter, and when utilised on real K-distributed clutter, they lose both detection performance and the constant false-alarm rate property. The optimum K-distributed detector of Gini \textit{et al.}~\cite{Gini1998} achieves the best possible performance, but requires the speckle covariance and texture statistics to be known, which is not the case in practice. Whilst practical approaches to detection in heterogeneous clutter exist, none exploit the optimum structure of this detector, and hence the aim of this thesis is to develop a practical implementation of this detector. This involves benchmarking the optimum detector against the Gaussian detectors, implementing it with estimated statistics, and investigating the possibility of integrating existing practical detectors with the optimum detector. So far, a K-distributed clutter generator has been developed and validated and preliminary results show the superior detection performance of the optimum detector compared to classical Gaussian detectors. This report mainly presents a review of the relevant literature, which provides the background needed for the remainder of the research project.